\begin{abstract}
Healthcare systems are increasingly leveraging large language models (LLMs) to aid in clinical decision-making, medical documentation, and other healthcare applications.
However, these models may exhibit stereotypical biases based on patient demographics and social determinants of health (SDOH), which 
impact patient care and clinical outcomes.
In this study, we tested whether AI systems are perpetuating existing disparities in healthcare by identifying and quantifying potential stereotypical biases in LLM-generated patient explanations.
We examined how LLMs from Meta, Anthropic, and OpenAI generate radiology report explanations tailored to patients of different genders, races/ethnicities,
socioeconomic statuses, and insurance types.
We used radiology reports from the MIMIC-IV corpus, applying four readability metrics: Flesch-Kincaid Grade Level, Simple Measure of Gobbledygook, Gunning Fog Index, and Coleman–Liau Index.
Across latest LLMs, privileged-group outputs consistently showed higher readability grades across all four indices.
Advanced models demonstrated large, statistically significant differences ranging from $\Delta \approx 0.8$ to $2.0$ grade levels, with effect sizes often exceeding $d>1.5$.
GPT-4-turbo exhibited the largest gaps across metrics (FKG: $\Delta = 1.669$, $d = 2.443$; GUNN: $\Delta = 1.976$, $d = 2.536$; all $p < 10^{-6}$).
Early-generation models showed minimal or nonsignificant differences.
Sentence Encoder Association Test (SEAT) analysis confirmed these patterns at the embedding level, with modern models showing significant associations of privileged-group embeddings with expert-like language (GPT-4-turbo: $d = -0.629$, $p < 10^{-6}$).
Demographic parity analysis revealed systematic disparities: privileged groups were overrepresented in high-complexity outputs (top-20\%: 22.09\% vs. 18.21\%), distributional differences were substantial (Wasserstein distance $W = 0.6073$), and counterfactual analysis showed demographic attributes alone caused higher complexity in 64.8\% of matched pairs.
These findings demonstrate that as LLMs become more advanced, their outputs exhibit greater linguistic complexity for privileged demographic groups, raising concerns for fairness and accessibility in patient communication.
\end{abstract}